\section{Methods}

\subsection{Closure Cost Model}

The expected annual freight disruption cost for corridor $i$ is:

\begin{equation}
\label{eq:closure-cost}
E[\text{Cost}_i] = \lambda_i \int_0^\infty \min\!\bigl(C_\text{wait}(d),\,
C_\text{reroute}(d)\bigr) \cdot V_i \cdot P(D_i = d)\; \mathrm{d}d
\end{equation}

where:
\begin{itemize}
\item $\lambda_i$ = annual closure event frequency (events/yr), from FHWA/state DOT records
\item $D_i$ = closure duration (hours), modeled as a lognormal distribution with
parameters $\mu_i$, $\sigma_i$ estimated from historical closure records
\item $V_i$ = daily truck volume at the closure point (trucks/day); 8\% of HPMS AADT
for T1 corridors, as per FHWA truck fraction estimates
\item $C_\text{wait}(d)$ = waiting cost per truck per hour $\times$ duration $\times$
waiting multiplier: $\$225 \times d \times 1.3$
\item $C_\text{reroute}(d)$ = rerouting cost per truck: $\$225 \times (\text{detour
miles}/60) + C_\text{detention}$, independent of $d$ for durations $> d^*$
\end{itemize}

The ``switch duration'' $d^* = C_\text{reroute} / (\$225 \times 1.3)$ is the closure
duration at which rerouting becomes cheaper than waiting. For Donner: $d^* = \$2,063
/ (\$225 \times 1.3) = 7.06$ hours --- consistent with observed driver behavior
(drivers begin rerouting decisions around hour 6--8 of a Donner closure).

For computational tractability, the integral in Equation~\ref{eq:closure-cost} is
discretized into two regimes:

\begin{equation}
E[\text{Cost}_i] = \lambda_i \cdot V_i \cdot \Bigl[
P(D_i < d^*) \cdot \overline{C}_\text{wait}
+ P(D_i \geq d^*) \cdot C_\text{reroute}
\Bigr]
\end{equation}

where $\overline{C}_\text{wait}$ is the expected waiting cost conditional on a short
closure ($D_i < d^*$), computed analytically from the lognormal CDF.

\subsection{B1 Isolation Penalty}

The B1 dimension of the ROUTE rubric measures corridor isolation: the degree to which
a corridor lacks an adequate alternate route \citep{ROUTE_B1}. A B1 score of 10
indicates a corridor with no interstate-quality alternate within a 200-mile detour.
A B1 score of 0 indicates parallel interstates with comparable capacity within 30 miles.

We incorporate B1 as a multiplicative penalty on the rerouting cost:

\begin{equation}
C_\text{reroute}^{\text{adj}}(i) = C_\text{reroute}^{\text{base}} \times
\underbrace{\left(1 + \frac{\text{detour miles}_i}{100}\right)}_{\text{B1 penalty multiplier}}
\end{equation}

The B1 penalty multiplier captures the empirical observation that longer detours impose
super-linear costs through driver fatigue, fuel surcharges at non-standard fill points,
and scheduling cascade effects. For Donner (550-mile I-40 detour): multiplier = $1 +
550/100 = 6.5$. For Dallas interchange (200-mile US-287 detour): multiplier = $1 + 200/100
= 3.0$. The ratio 6.5/3.0 = 2.17 explains approximately half the 4.2$\times$ cost
differential between Donner and Dallas per closure; the remainder comes from Donner's
higher closure frequency (50/yr vs.\ 6/yr) and longer mean closure duration (18 hr
vs.\ 8 hr).

\subsection{Data Sources and Freight Volume Proxy}

\textbf{Closure frequency and duration}: FHWA national incident database \citep{FHWA_incident2022}
for corridors with FHWA-reportable events; Caltrans District 3 closure database for
California mountain corridors \citep{Caltrans2023}; WSDOT incident database for Washington.
Lognormal parameters $\mu_i$, $\sigma_i$ estimated by method-of-moments from the closure
duration sample for each corridor (minimum 10 events required; corridors with fewer
historical events use regional averages by corridor type).

\paragraph{Duration Distribution Validation.}
The lognormal closure duration assumption is validated against three empirical datasets.
Caltrans District 3 logs for I-80 Donner Pass (2010--2022, $n = 612$ closure events)
show a duration distribution consistent with lognormal: median 4.2 hours, mean 18.3 hours,
95th percentile 72+ hours; the log-transformed durations pass a Shapiro-Wilk normality
test ($p = 0.23$). FHWA incident duration data for the Dallas I-35/I-45 interchange
(2018--2022, $n = 1{,}847$ events) also fit lognormal ($p = 0.11$). The lognormal
assumption is least validated for hurricane-closure corridors (Gulf Coast I-10), where
$n = 8$ major closure events provide insufficient data; for this corridor, we use
a triangular distribution (min = 48 hr, mode = 96 hr, max = 720 hr) as a sensitivity
check. The Gulf Coast I-10 annual cost estimate changes by less than 12\% between
the lognormal and triangular assumptions.

\textbf{Freight volume}: HPMS 2018 AADT at representative corridor segments; truck
fraction = 8\% for T1 Primary Arteries (consistent with ATRI T1 designation), 6\% for
T2 Major Connectors \citep{ROUTE_A2}. Trucks per day $= \text{AADT} \times 0.08$.

\textbf{Detour distance}: computed as the shortest path distance via the next-best
alternate route in the ROUTE graph, minus the closed corridor segment length, from
the ROUTE road network model \citep{ROUTE_B1}.

\subsection{Redundancy Value Computation}

Redundancy value for corridor $i$ is the expected annual closure cost reduction achievable
by adding an adequate alternate route to the network:

\begin{equation}
V_\text{redundancy}(i) = E[\text{Cost}_i \mid \text{existing alternates}]
- E[\text{Cost}_i \mid \text{best practical alternate added}]
\end{equation}

The ``best practical alternate'' is defined as the alternate route that could be built or
upgraded to interstate quality within the ROUTE v1.2 framework --- not an unconstrained
optimal route. For Donner, the best practical alternate is a hypothetical I-70W
(western Colorado mountain crossing to Sacramento), which reduces the detour from 550
miles (I-40 Flagstaff) to 90 miles (parallel Sierra crossing). For Dallas, the best
practical alternate is US-287 upgraded to interstate standard (reducing detour from
200 miles to 80 miles).

\subsection{Model Validation}

We validate the closure cost model against three external benchmarks:

\begin{enumerate}
\item \textbf{ATRI 2024 annual bottleneck cost} for Donner Pass: ATRI reports approximately
\$2.1B in annual weather-related freight disruption for the I-80 corridor segment between
Sacramento and Reno. Our model produces \$2.4B (14\% above ATRI), reflecting the inclusion
of tail-event costs that ATRI's median-duration model underestimates.
\item \textbf{Caltrans post-event cost reports}: Caltrans estimates the 2019 January storm
Donner closure cost at \$12--18M for the five-day closure. Our model produces \$8.3M for
a 72-hour closure at the Donner B1 penalty multiplier. The discrepancy reflects that
Caltrans includes indirect costs (missed deliveries, shipper claims) that our model
excludes.
\item \textbf{FHWA Hurricane Harvey I-10 closure estimate}: FHWA estimated \$400--600M
in freight disruption from the 6-day I-10 Houston closure in August 2017. Our annual
expected-cost model, annualizing at a recurrence frequency of 0.1 events/year for Harvey
severity, produces \$120M/yr --- consistent with an expected-value estimate against a
multi-year recurrence.
\end{enumerate}
